\section{Qu’est-ce que la gestion de projet ?}

\begin{frame}[t]{Qu’est-ce qu’un projet ?}
    \textbf{Projet :} processus permettant d’atteindre un objectif fixé.

    \vspace{1.4cm}

    \begin{center}
    \begin{tikzpicture}
        \node[image] (idea) {\includegraphics[width=1cm]{icon-idea}\\Idée};
        \node[image] (project) at (8cm, 0) {\includegraphics[width=1cm]{icon-target}\\Objectif};

        \draw[->, shortarrow] (idea) -- node[below] {\textit{Projet}} (project);
    \end{tikzpicture}
    \end{center}

    \note{note test}
\end{frame}

\begin{frame}[t]{Comment gérer un projet ?}
    \textbf{Gestion de projet :} outils et méthodes permettant d’organiser un projet.

    \vfill

    \begin{tikzpicture}
        \node[image] (project) {\includegraphics[width=1cm]{icon-project}\\Projet};
        \node[image] (team) at (\textwidth - 2cm, 0) {\includegraphics[width=1cm]{icon-team}\\Équipe};

        {[visible on=<2->]
            \node[image, right=2cm of project] (task2) {\includegraphics[width=1cm]{icon-task}};
            \node[image, above=.5cm of task2] (task1) {\includegraphics[width=1cm]{icon-task}};
            \node[image, below=.5cm of task2] (task3) {\includegraphics[width=1cm]{icon-task}\\Tâches};

            \draw[->, shortarrow] (project) edge[out=20, in=200] (task1);
            \draw[->, shortarrow] (project) -- (task2);
            \draw[->, shortarrow] (project) edge[out=-20, in=160] (task3);
        }

        {[visible on=<3->]
            \node[image, left=2cm of team] (role2) {\includegraphics[width=1cm]{icon-man}};
            \node[image, above=.5cm of role2] (role1) {\includegraphics[width=1cm]{icon-woman}};
            \node[image, below=.5cm of role2] (role3) {\includegraphics[width=1cm]{icon-woman}\\Rôles};

            \draw[->, shortarrow] (team) edge[out=160, in=-20] (role1);
            \draw[->, shortarrow] (team) -- (role2);
            \draw[->, shortarrow] (team) edge[out=200, in=20] (role3);
        }

        {[visible on=<4->]
            \draw[<->, shortarrow] (task1) -- (role1);
            \draw[<->, shortarrow] (task1) -- (role3);
            \draw[<->, shortarrow] (task2) -- (role1);
            \draw[<->, shortarrow] (task3) -- (role2);
        }
    \end{tikzpicture}
\end{frame}

\newcommand{\imgcaption}[2]{%
    \begin{minipage}[t]{0.135\linewidth}%
        \centering
        \includegraphics[width=\linewidth]{#1}\\%
        {\scriptsize #2}%
    \end{minipage}%
}

\begin{frame}{Pourquoi utiliser une méthode ?}
\makebox[\linewidth][c]{%
    \begin{minipage}{\dimexpr\textwidth+4em\relax}
        \imgcaption{images/arbre-1}{Comment le client a exprimé ses besoins}%
        \hfill%
        \imgcaption{images/arbre-2}{Comment le chef de projet les a compris}%
        \hfill%
        \imgcaption{images/arbre-3}{Ce que les ingénieurs ont conçu}%
        \hfill%
        \imgcaption{images/arbre-4}{Ce que les programmeurs ont codé}%
        \hfill%
        \imgcaption{images/arbre-5}{Comment le projet a été vendu}%
        \hfill%
        \imgcaption{images/arbre-6}{Ce qu’à reçu le client}%
        \hfill%
        \imgcaption{images/arbre-7}{Ce que voulait vraiment le client}%
    \end{minipage}%
}
\end{frame}
